%------------------------------------------------------------------------------%

\usepackage{calculo_varias_variaveis-1}

\title {Fórmulas de De Moivre}

%------------------------------------------------------------------------------%
\begin{document}

\addtitle

%------------------------------------------------------------------------------%
\section{Potencias de Números Complexos}

%------------------------------------------------------------------------------%
\begin{frame}{Potencias}
  \begin{align*}
    \onslide<+->{z^n}
    \onslide<+->{& = zz\cdots z \\[5mm]}
    \onslide<+->{& = \rho\rho\cdots\rho}
    \onslide<+->{    \big[\cos(\varphi\varphi\cdots\varphi)
      + i\sin(\varphi\varphi\cdots\varphi)
      \big] \\[5mm]}
    \onslide<+->{& = \rho^n\big[\cos(n\varphi) + i\sin(n\varphi)\big]}
  \end{align*}
\end{frame}

%------------------------------------------------------------------------------%
\begin{frame}{Primeira Fórmula de De Moivre}
  Se
  \[
    z = \rho \big[\cos(\varphi) + i\sin(\varphi)\big]
  \]
  e $n$ um número natual

  \pause
  \vspace{\baselineskip}
  Primeira Fórmula de De Moivre
  \[
    z^n = \rho^n \big[\cos(n\varphi) + i\sin(n\varphi)\big]
  \]
\end{frame}

%------------------------------------------------------------------------------%
\addtocounter{exemplo}{1}
\begin{frame}{Exemplo \theexemplo}
  Calcule
  \[
    \left[2\left(\cos\left(\frac{\pi}{6}\right) + i\sin\left(\frac{\pi}{6}\right)\right)\right]^4
  \]
\end{frame}

\begin{frame}{Exemplo \theexemplo}
  \[
    z = 2\left(\cos\left(\frac{\pi}{6}\right) + i\sin\left(\frac{\pi}{6}\right)\right)
  \]

  \[
    \pause
    \rho = 2
    \qquad\qquad
    \pause
    \varphi = \frac{\pi}{6}
  \]

  \[
    \pause
    z^4
    \pause
    = \rho^4 \big[\cos(4\varphi) + i\sin(4\varphi)\big]
  \]
  \[
    \hphantom{z^4}
    \pause
    = 2^4 \left[\cos\left(4\frac{\pi}{6}\right) + i\sin\left(4\frac{\pi}{6}\right)\right]
    \pause
    = 16  \left[\cos\left(\frac{2\pi}{3}\right) + i\sin\left(\frac{2\pi}{3}\right)\right]
  \]
\end{frame}

%------------------------------------------------------------------------------%
\addtocounter{exemplo}{1}
\begin{frame}{Exemplo \theexemplo}
  Calcule \(\left(1+\sqrt{3}i\right)^6\)
\end{frame}

\begin{frame}{Exemplo \theexemplo}
  \[
    z = 1+\sqrt{3}i
    \qquad\qquad
    \pause
    a = \Re(z) = 1
    \qquad\qquad
    \pause
    b = \Im(z) = \sqrt{3}
  \]

  \[
    \pause
    \rho
    \pause
    = \sqrt{a^2 + b^2}
    \pause
    = \sqrt{1 + 3}
    \pause
    = \sqrt{4}
    \pause
    = 2
  \]

  \[
    \pause
    \sin(\varphi)
    \pause
    = \frac{b}{\rho}
    \pause
    = \frac{\sqrt{3}}{2}
    \pause
    \qquad\qquad
    \cos(\varphi)
    \pause
    = \frac{a}{\rho}
    \pause
    = \frac{1}{2}
  \]

  \[
    \pause
    \varphi
    \pause
    = \ang{60}
    \pause
    = \frac{\pi}{3} \si{\radian}
  \]

  \[
    \pause
    z
    \pause
    = \rho\left(\cos(\varphi)+i\sin(\varphi)\right)
    \pause
    = 2\left(\cos\left(\frac{\pi}{3}\right)+i\sin\left(\frac{\pi}{3}\right)\right)
  \]
\end{frame}

\begin{frame}{Exemplo \theexemplo}
  \[
    z = 1+\sqrt{3}i
    \pause
    = 2\left(\cos\left(\frac{\pi}{3}\right)+i\sin\left(\frac{\pi}{3}\right)\right)
  \]

  \vspace{0\baselineskip}
  \[
    \pause
    z^6
    \pause
    = \rho^6 \big[\cos(6\varphi) + i\sin(6\varphi)\big]
    \]
    \[
    \hphantom{z^6}
    \pause
    = 2^6 \left[\cos\left(6\frac{\pi}{3}\right) + i\sin\left(6\frac{\pi}{3}\right)\right]
    \pause
    = 64  \left[\cos\left(2\pi\right) + i\sin\left(2\pi\right)\right]
    \pause
    = 64
  \]
\end{frame}

%------------------------------------------------------------------------------%
\section{Raizes de Números Complexos}

%------------------------------------------------------------------------------%
\begin{frame}{Raizes}
  Se
  \(
    z = \rho \big[\cos(\varphi) + i\sin(\varphi)\big]
  \)
  e $n>1$ é um número natual

  \pause
  \vspace{\baselineskip}
  A raiz $n$-ésima de $z$ é um número
  \[
    u = r \big[\cos(\alpha) + i\sin(\alpha)\big]
  \]
  \pause
  tal que
  \[
    u^n = z
  \]
  \[
    \pause
    r^n \big[\cos(n\alpha) + i\sin(n\alpha)\big]
    = \rho \big[\cos(\varphi) + i\sin(\varphi)\big]
  \]
\end{frame}

%------------------------------------------------------------------------------%
\begin{frame}{Raizes}
\begin{columns}
\column{0.5\linewidth}
  Encontrar $r$ e $\alpha$
  \pause
   tais que
  \begin{gather*}
    r^n = \rho \\
    \cos(n\alpha) = \cos(\varphi) \\
    \sin(n\alpha) = \sin(\varphi)
  \end{gather*}
  \pause
  ou seja
  \begin{gather*}
    r=\sqrt[n]{\rho} \\
    n\alpha = \varphi + 2k\pi
  \end{gather*}
  \column{0.5\linewidth}
  \pause
  \[
    \alpha = \frac{\varphi + 2k\pi}{n}
  \]
  \pause
  apenas
  \[
    k = 0, 1, 2, \dots, n-1
  \]
  fornecem soluções distintas
\end{columns}
\end{frame}

%------------------------------------------------------------------------------%
\begin{frame}{Sengunda Formula de De Moire}
  Se
  \(
    z = \rho \big[\cos(\varphi) + i\sin(\varphi)\big]
  \)
  e $n>1$ é um número natual

  \vspace{\baselineskip}
  As raizes $n$-ésimas de $z$ são
  \[
    u_k = \sqrt[n]{\rho}
    \left[
         \cos\left(\frac{\varphi+2k\pi}{n}\right)
      + i\sin\left(\frac{\varphi+2k\pi}{n}\right)
    \right]
  \]
  com \(k=0, 1, 2, \dots, n-1\)
\end{frame}

%------------------------------------------------------------------------------%
\addtocounter{exemplo}{1}
\begin{frame}{Exemplo \theexemplo}
  Encontre as raizes cúbicas de 1
\end{frame}

\begin{frame}{Exemplo \theexemplo}
  \[
    z = 1
    \pause
    = 1 + 0i
    \pause
    = 1\big[\cos(0) + i\sin(0)\big]
  \]

  \[
    \pause
    \rho = 1
    \qquad\qquad
    \pause
    \varphi = 0
  \]

  \[
    \pause
    u_k = \sqrt[3]{\rho}
    \left[
       \cos\left(\frac{\varphi+2k\pi}{3}\right)
    + i\sin\left(\frac{\varphi+2k\pi}{3}\right)
    \right]
    \qquad\qquad
    k=0, 1, 2
  \]
  \[
    \pause
    \hphantom{u_k}
    = \sqrt[3]{1}
    \left[
    \cos\left(\frac{0+2k\pi}{3}\right)
    + i\sin\left(\frac{0+2k\pi}{3}\right)
    \right]
  \]
  \[
    \pause
    \hphantom{u_k}
    =  \cos\left(\frac{2k\pi}{3}\right)
    + i\sin\left(\frac{2k\pi}{3}\right)
  \]
\end{frame}

\begin{frame}{Exemplo \theexemplo}
\vspace{-1.5\baselineskip}
\begin{columns}[t]
\column{0.5\linewidth}
  \begin{align*}
  \onslide<+->{u_0}
  \onslide<+->{& = \cos\left(\frac{2\times 0 \pi}{3}\right) + i\sin\left(\frac{2\times 0 \pi}{3}\right) \\}
  \onslide<+->{& = \cos(0) + i\sin(0) \\}
  \onslide<+->{& = 1 \\[5mm]}
  \onslide<+->{u_1}
  \onslide<+->{& = \cos\left(\frac{2\times 1 \pi}{3}\right) + i\sin\left(\frac{2\times 1 \pi}{3}\right) \\}
  \onslide<+->{& = \cos\left(\frac{2\pi}{3}\right) + i\sin\left(\frac{2\pi}{3}\right) \\}
  \onslide<+->{& = -\frac{1}{2} + \frac{\sqrt{3}}{2}i}
\end{align*}
\column{0.5\linewidth}
\begin{align*}
  \onslide<+->{u_2}
  \onslide<+->{& = \cos\left(\frac{2\times 2 \pi}{3}\right) + i\sin\left(\frac{2\times 2 \pi}{3}\right) \\}
  \onslide<+->{& = \cos\left(\frac{4\pi}{3}\right) + i\sin\left(\frac{4\pi}{3}\right) \\}
  \onslide<+->{& = -\frac{1}{2} - \frac{\sqrt{3}}{2}i}
\end{align*}
\end{columns}
\end{frame}

%------------------------------------------------------------------------------%
\section{Lista Mínima}

%------------------------------------------------------------------------------%
\begin{frame}{Lista Mínima}
%%  Cálculo Vol. 2 do Thomas 12\textsuperscript{a} ed. -- Seção 14.8
%%  \begin{enumerate}
%%    \item Estudar o texto da seção
%%    \item Resolver os exercícios: 3, 5, 7, 12, 16, 18, 24 e 30
%%  \end{enumerate}

  \vfill
  Atenção: A prova é baseada no livro, não nas apresentações
\end{frame}

%------------------------------------------------------------------------------%
\end{document}
%------------------------------------------------------------------------------%
